\documentclass[11pt,a4paper]{article}
\usepackage[T2A]{fontenc}
\usepackage[utf8]{inputenc}
\usepackage[russian]{babel}
\usepackage[margin=21mm]{geometry}
\usepackage{amsmath,amssymb,longtable,booktabs,array}
\usepackage[unicode,hidelinks]{hyperref}
\setlength{\parindent}{0pt}
\setlength{\parskip}{5pt}
\setlength{\emergencystretch}{3em}
\begin{document}
\title{Фальсификация нагрузки SMF и независимая проверка через SCP  воспроизводимый эксперимент на Open5GS}
\author{}\date{}\maketitle

\section*{Аннотация}

Проверена причинная цепочка фальсификации нагрузки SMF в изолированном стенде Open5GS v2.8.0. Исходная конфигурация назначала все 900 проверенных сессий первому зарегистрированному SMF. После введения свежего discovery и взвешенного выбора выполнено 16 пригодных парных сравнений честного и заниженного отчёта. Средние парные эффекты в доступных сочетаниях политики и нагрузки составили от 1.87 до 8.50 процентного пункта. Независимое наблюдение SCP позволяет сравнивать отчёт с числом обслуживаемых контекстов без использования метрик SMF в признаках. Разделены время доставки отчёта в NRF, его возраст при выборе и ошибки детекторов на отдельных тестовых запусках. Вывод ограничен исследовательской политикой и пилотным объёмом; полное наблюдение и доверенная шкала ёмкости являются существенными предпосылками. 

\section*{Постановка}

Исследуется влияние недостоверного отчёта SMF о нагрузке на распределение новых PDU-сессий. Атакующий контролирует одну из трёх SMF и изменяет собственные отчёты, сохраняя возможность обслуживать сессии. NRF, AMF, SCP и генератор UE считаются доверенными. Атакующий не изменяет журналы SCP и не получает доступ к результату работы детектора. Стенд изолирован; внешняя сеть оператора не используется.

Первичная величина — изменение доли успешных созданий контекста у SMF3:

\[
\Delta S=S_3^{attack}-S_3^{honest},\qquad
S_i={N_i^{successful\ CreateSMContext}\over\sum_jN_j^{successful\ CreateSMContext}}.
\]

Это показатель смещения назначения сессий. Сам по себе он не доказывает перехват пользовательского трафика, отказ в обслуживании или нарушение задержек приложения. Для таких выводов потребовались бы отдельные измерения пользовательской плоскости.

\section*{Аудит реализации и стенд}

Использованы Open5GS v2.8.0 (\texttt{157f611a530e292e40ec50f9d 23f0ef5d4fcd6a6}) и UERANSIM v3.2.7 (\texttt{1d1e154f869260b5e98f69058 27b1bd9b8663afc}). В стенд входят три SMF, по одному NRF, SCP, AMF, UPF и gNB, вспомогательные AUSF/UDM/UDR/PCF/NSSF/BSF, MongoDB и Prometheus. Все SMF имеют одинаковые DNN, S-NSSAI и TAI. Пулы адресов SMF раздельны. Функции работают в одном Linux-контейнере; общие монотонные часы позволяют сопоставлять события без сетевой синхронизации времени.

Проверка исходников уточнила исходную гипотезу: SMF v2.8.0 уже вычисляет долю занятости пула перед обработкой SBI-таймеров. Поэтому утверждение, что её отчёт всегда равен нулю, неверно. Однако обработчик NRF не присваивает новое значение из PATCH \texttt{/load}, а штатный выбор первого подходящего NF не реализует исследуемую взвешенную политику. Сведения относятся к закреплённой версии, а не ко всем версиям Open5GS. \href{https://github.com/open5gs/open5gs/tree/v2.8.0}{Исходники Open5GS v2.8.0}.

Функциональные изменения:

\noindent 1. SMF считает успешные SBI-контексты и перед heartbeat формирует операционную нагрузку. Честная и атакующая SMF используют одну сборку с разными параметрами запуска.\par
\noindent 2. NRF применяет целочисленный \texttt{/load} в диапазоне 0–100. Отдельная проверка HTTP/2 сопоставляет исходный и исправленный NRF, включая шесть некорректных значений.\par
\noindent 3. AMF выполняет свежий discovery для нового контекста и выбирает SMF среди подходящих результатов по описанной ниже политике.\par
\noindent 4. SCP регистрирует успешные создания, освобождения и уведомления RELEASED; последующий анализ реконструирует число контекстов независимо от метрик SMF.\par
\noindent 5. Номер версии отчёта передаётся диагностическими HTTP-заголовками. Он нужен для сопоставления времён и не входит в признаки детектора.\par

В UERANSIM исправлен экспериментально обнаруженный конфликт CLI с приостановкой потоков: операции установления и освобождения передаются очередью в NAS. Генератор ждёт фактическое подтверждение операции в журнале UE; ответ CLI сам по себе не считается успешной PDU-сессией. \href{https://github.com/aligungr/UERANSIM/wiki/Usage}{Документация UERANSIM}.

После одновременной потери эмулируемой радиосвязи с множеством UE тайм-аут heartbeat радиосимулятора увеличен с 2 до 15 секунд. Это параметр UERANSIM, не секундный heartbeat нагрузки SMF. Изменение отмечено в manifest новых запусков; неудачная трасса сохранена и исключена. Успешные более ранние опыты остаются отдельной версией конфигурации; влияние этого изменения является ограничением сопоставимости серий.

Оставшиеся опыты перенесены на четыре независимые копии стенда с раздельными сетевыми пространствами, базами данных и непересекающимися наборами по четыре CPU. Измерительные бинарные файлы скопированы из одной сборки. Лимит времени учитывается суммарно по активным прогонам, а не отдельно для каждой копии. Ранее полученные последовательные опыты явно сохраняют свой режим в метаданных. Общий хост и различие режимов ограничивают интерпретацию точных задержек и межсерийных сравнений.

\section*{Нагрузка и выбор SMF}

Для основной серии абсолютная экспериментальная ёмкость каждого SMF равна 100 сессиям:

\[
L_i^{ref}=\min(100,100N_i^{active}/C_i^{abs}),\qquad
L_i^{rep}=\lfloor(1-\alpha_i)L_i^{ref}\rfloor.
\]

\texttt{C\_abs} задаёт операционную шкалу и не является измеренным аппаратным пределом. Параметр \texttt{capacity} NFProfile — отдельный относительный вес. В основной серии он одинаков и равен 100. Счётчик SMF увеличивается при успешном ответе CreateSMContext и уменьшается при удалении контекста; завершение процедуры на UE проверяется дополнительно.

Сначала применяются фильтры discovery, затем остаются NF с наименьшим численным \texttt{priority}. Для них задаётся вес

\[
w_i=capacity_i(100-L_i^{rep}).
\]

Проверяются SWRR и выбор с вероятностью \(w_i/\sum_jw_j\). При нулевых весах предусмотрен переход к относительной ёмкости, а затем к равномерному распределению. Основной режим не рассчитан на все возможные сочетания NF-level и service-level параметров. Стандарт задаёт смысл \texttt{priority}, \texttt{capacity}, \texttt{load} и процедуры NRF, но приведённая формула и выбор SWRR являются экспериментальной политикой. \href{https://www.etsi.org/deliver/etsi_ts/129500_129599/129510/18.11.00_60/ts_129510v181100p.pdf}{3GPP TS 29.510 / ETSI V18.11.0}.

\section*{Аналитическая модель}

При одинаковой ёмкости, постоянном общем числе сессий, равномерном выборе освобождаемой активной сессии и стационарном режиме доля занятых контекстов равна доле новых назначений. Для одного атакующего узла обозначим её \(s\), а общую нормированную нагрузку — \(\rho=N/(3C^{abs})\). Тогда нагрузка атакующего равна \(3\rho s\), а каждого честного узла — \(3\rho(1-s)/2\). Из нормировки весов следует

\[
3\rho\alpha s^2+(3-3\rho\alpha)s-1=0,
\qquad
s={2\over3-a+\sqrt{(3-a)^2+4a}},\quad a=3\rho\alpha.
\]

Честная симметричная точка — \(1/3\). При \(\alpha=0.5\) модель предсказывает прирост около 2.29, 5.96 и 8.52 процентного пункта для \(\rho=0.2,0.5,0.7\). Это расчёт, а не измерение. Он применим вне насыщения, при \(3\rho s<1\), и не учитывает целочисленное округление, heartbeat, конечное время разогрева и случайность. Согласие пилотных точек с моделью не заменяет проверки стационарности.

Один и тот же отчёт может быть сформирован честным малонагруженным узлом или более нагруженным узлом, который занижает значение. Поэтому одного численного отчёта без дополнительных предпосылок недостаточно для идентификации атаки. Межузловое сравнение в NRF-only является эвристикой; независимое наблюдение SCP добавляет информацию о фактически созданных контекстах.

\section*{Протокол и разделение данных}

Первоначальный расширенный протокол предусматривал десять seeds и 3000 измеряемых назначений в каждом опыте. Он не объявляется выполненным. С учётом разрешённых четырёх часов работы стенда зафиксирована пилотная серия: три отдельных честных калибровочных запуска, затем парные опыты с \(\alpha=0\) и \(0.5\), двумя политиками и тремя уровнями нагрузки. План включает три независимых seed; фактическое число завершённых повторов приведено в результатах.

В основном опыте после начального заполнения исключаются 200 замен сессий, затем измеряются 400. В калибровке — 100 и 200. В дополнительных опытах устойчивости до их запуска принят объём 100+200. Генератор поддерживает постоянное число активных UE: случайно выбирает активную сессию для освобождения и свободный UE для создания новой. Парные опыты используют одинаковую последовательность операций; её хеш проверяется. Время поступления не идентично, поскольку генератор ждёт подтверждений. UUID узлов между запусками меняются, поэтому одинаковый seed случайного выбора не означает одинаковые назначения физическим SMF.

Независимая единица повторения — целый запуск, а не сессия и не десятисекундное окно. Средний парный эффект рассчитывается по запускам. При наличии не менее трёх пар приводится 95\%-й percentile bootstrap-интервал по 10 000 выборок пар; при трёх повторах его следует считать разведочным. Окна не трактуются как независимые повторы для проверки значимости.

В первичный средний эффект входят только пары с одинаковым тайм-аутом радиосимулятора и одинаковым режимом ресурсов — последовательным либо параллельным с четырьмя выделенными CPU. Переходные пары сохраняются в CSV с признаком \texttt{comparable=false}. Если опыт повторён для согласования конфигураций, выбирается наиболее поздняя подходящая пара; правило не зависит от величины эффекта. Все исходные прогоны остаются в аудите.

Из анализа исключаются технически незавершённые опыты, расхождения количества успешных установлений с UE, повторные события создания, несовпадение итоговых счётчиков SCP и SMF, неверное число измеряемых назначений и несоответствие версии/значения нагрузки AMF отправленному отчёту. Неудачные журналы сохраняются. Промежуточная разница счётчиков из-за порядка асинхронных событий отдельно отражается в аудите.

\section*{Независимые признаки и детекторы}

SCP сопоставляет URI созданного контекста, целевую SMF и callback. Успешный CreateSMContext увеличивает счётчик; успешный Release или уведомление RELEASED удаляет тот же контекст, предотвращая повторное уменьшение. Детектор использует

\[
\widehat L_i^{SCP}=\min(100,100\widehat N_i^{SCP}/C_i^{abs}),\qquad
D_i=|L_i^{rep}-\widehat L_i^{SCP}|.
\]

\texttt{sessionnbr}, внутренний reference load и внутренний active count SMF используются только для контроля измерений. Изменение этих контрольных данных в тестовой копии журнала не меняет признаки. Известная операционная ёмкость — доверенные метаданные стенда; возможность получить её независимо в рабочей сети является отдельным условием применимости.

В каждом полном десятисекундном окне, содержащем не менее пяти назначений, вычисляются шесть простых оценок:

{\footnotesize\setlength{\tabcolsep}{3pt}\renewcommand{\arraystretch}{1.25}
\begin{longtable}{p{8.16cm}p{8.16cm}}
\toprule
Оценка & Данные и смысл \\
\midrule\endhead
NRF-only & Максимальное отклонение заявленной нагрузки SMF3 от медианы двух других SMF \\
SCP-only & Отклонение доли новых назначений SMF3 от 1/3, в процентных пунктах \\
Residual & Максимум \(D_3\) по обновлениям NRF в окне \\
Residual quantized & То же сравнение после целочисленного округления наблюдаемой нагрузки вниз, как в SMF \\
NRF+SCP & Максимум residual и отклонения относительной capacity от доверенной конфигурации 100 \\
Full peer & Дополнительно разность межузловых контрастов заявленной и наблюдаемой нагрузки \\
\bottomrule\end{longtable}}

Порог каждой оценки — 99-й эмпирический процентиль только честной калибровки. Решение — строгое превышение порога. На тестовой части отдельно измеряется реальный FPR: настройка по 99-му процентилю не гарантирует FPR=1\%. Pd основной атаки относится только к постоянному занижению \(\alpha=0.5\); результаты адаптивных и меняющихся во времени воздействий приводятся по отдельным кампаниям. Число ложноположительных окон в час не равно числу отдельных эпизодов тревоги.

Три калибровочных запуска используют SWRR. Применение фиксированного порога к вероятностному выбору проверяет переносимость между политиками; это не отдельная калибровка вероятностной политики. Поэтому FPR и Pd дополнительно разделены по политике, а сравнение при одинаковом достигнутом FPR=1\% не заявляется.

Проводится абляция доступных источников данных, а не обучение нейросети. Добавление ML без преимущества над простым residual в этой постановке не обосновано.

Вариант с учётом округления добавлен до выполнения неоднородного честного контроля. При ёмкости 100 нагрузка всегда целая, поэтому калибровка не учит порог ошибке округления. При ёмкости 70 исходный residual может реагировать на честную дробную нагрузку; вариант quantized проверяет эту конкретную причину ложных тревог без изменения исходного детектора.

\section*{Дополнительные воздействия и контроль}

Атакующий с ограничением остатка уменьшает отчёт на \(\min(\alpha L^{ref},b)\), где \(b\in\{0,1,2,4,8\}\) процентных пунктов. Это адаптация к собственной нагрузке и заданному бюджету искажения, без доступа к оценке детектора. \(b=0\) является нулевым воздействием и проверяет согласованность результатов. Отдельно предусмотрены завышение относительной capacity до 150, его сочетание с ограничением остатка, линейное нарастание и периодическое включение/выключение занижения. Для этих опытов используется отдельный честный контрфактический запуск с тем же seed.

Если ограничение активно и отчёт уменьшается на постоянное число пунктов b, ненасыщенная модель с равными capacity даёт долю \(s=(1+b/100)/(3+b/100)\). Соответствующие приросты для b=1,2,4,8 составляют приблизительно 0.22, 0.44, 0.88 и 1.73 процентного пункта. При 200 измеряемых назначениях дискретность доли составляет 0.5 пункта; такой пилот с одним seed не разрешает надёжно самые малые эффекты. Занижение нагрузки и завышение capacity необходимо различать: детектор, проверяющий только load, по построению не проверяет истинность относительного веса.

Честный неоднородный контроль задаёт SMF3 операционную ёмкость 70 вместо 100 и передаёт это значение наблюдателю как доверенные метаданные. Он проверяет реакцию на известную неоднородность шкалы нагрузки. Это не эксперимент по физическому замедлению CPU или задержке сети.

\section*{Время распространения}

Ключ корреляции — \texttt{(nfInstanceId, load\_version)}, а не совпадение численного load. Измеряются \(t_0\) вычисления, \(t_1\) отправки, \(t_2\) сохранения NRF и \(t_3\) использования AMF. Отдельно приводятся \(t_2-t_0\) и возраст отчёта \(t_3-t_0\). Последний включает ожидание очередного выбора и heartbeat раз в секунду; его нельзя целиком интерпретировать как сетевую задержку. Результаты одного контейнера не описывают географически распределённую сеть.

\section*{Ограничения и границы вывода}

Исследование проверяет механизм на модифицированной политике выбора. Оно не доказывает, что исходный Open5GS подвержен тому же смещению. Короткий разогрев и малое число независимых повторов ограничивают статистические выводы. Чистое полное наблюдение SCP, отсутствие потерь уведомлений и известная ёмкость делают residual сильным эталоном, но требуют отдельной проверки при рестартах, ретрансляциях, обходе SCP, потерях событий, изменении ёмкости и задержках. Такие сценарии нельзя считать выполненными по наличию соответствующего пункта в плане.

Уровень патчей — исследовательский: однородная локальная сеть, ограниченный набор NF, базовая проверка \texttt{/load}; не проводится заявление о полной реализации всех операций JSON Patch или пригодности для промышленного внедрения. Не выполнена отдельная юридическая оценка патентов. Приоритет идеи и научная новизна относительно всей литературы не установлены; экспериментальные результаты сами по себе этого не доказывают.



\section*{Фактически полученные результаты}


Срез данных: 2026-10-03T15:02:06+05:00. \textbf{Таблицы отражают сохранённые завершённые опыты.}

Аудит прошли 49 запусков из 49 технически завершённых. В них подтверждено 32340 успешных установлений PDU-сессий; после исключения заполнения, разогрева и калибровки — 16200 измеряемых назначений. Три исходных B0-прогона (900 установлений) учитываются отдельно.

Консервативный учёт суммарного времени активных опытов, включая отладку и неудачные попытки: 3 ч 56 мин 50 с при лимите 4 часа. Для параллельных копий времена складываются. В первичную оценку вошли 16 из 18 запланированных сопоставимых пар: для SWRR при 20\% одна пара пересекла изменение конфигурации, а для случайного выбора при 20\% третий парный повтор не запускался из-за лимита. Переходные результаты не скрываются.

\subsection*{Исходный B0}

{\footnotesize\setlength{\tabcolsep}{3pt}\renewcommand{\arraystretch}{1.25}
\begin{longtable}{p{3.94cm}p{3.94cm}p{3.94cm}p{3.94cm}}
\toprule
Порядок регистрации & SMF1 & SMF2 & SMF3 \\
\midrule\endhead
1\ensuremath{\to}2\ensuremath{\to}3 & 300 & 0 & 0 \\
2\ensuremath{\to}3\ensuremath{\to}1 & 0 & 300 & 0 \\
3\ensuremath{\to}1\ensuremath{\to}2 & 0 & 0 & 300 \\
\bottomrule\end{longtable}}

При перестановке порядка все назначения следовали за первым зарегистрированным SMF. Это подтверждает влияние порядка в проверенной конфигурации исходной версии и объясняет необходимость явной экспериментальной политики.

\subsection*{Эффект занижения нагрузки}

{\footnotesize\setlength{\tabcolsep}{3pt}\renewcommand{\arraystretch}{1.25}
\begin{longtable}{p{2.53cm}p{2.53cm}p{2.53cm}p{2.53cm}p{2.53cm}p{2.53cm}}
\toprule
Политика & Нагрузка & Пар & \ensuremath{\Delta}S, п.п. & 95\% bootstrap, п.п. & Теория, п.п. \\
\midrule\endhead
random & 20\% & 2 & 2.00 & — \ldots{} — & 2.29 \\
random & 50\% & 3 & 6.25 & 6.00 \ldots{} 6.50 & 5.96 \\
random & 70\% & 3 & 8.50 & 7.75 \ldots{} 9.25 & 8.52 \\
swrr & 20\% & 2 & 1.87 & — \ldots{} — & 2.29 \\
swrr & 50\% & 3 & 7.25 & 5.00 \ldots{} 8.75 & 5.96 \\
swrr & 70\% & 3 & 7.58 & 6.00 \ldots{} 8.50 & 8.52 \\
\bottomrule\end{longtable}}

Интервал не выводится при числе пар меньше трёх. Даже три пары дают только разведочную оценку; знак эффекта одного опыта не служит доказательством универсального преимущества атакующего. В контрольных отсчётах SMF зафиксировано 0 случаев достижения операционной ёмкости из 38714 отсчётов. Они включают заполнение и разогрев; это индикатор нарушения ненасыщенного приближения, а не отдельная оценка временной доли перегрузки.

\subsection*{Детекторы на отложенных тестовых запусках}

{\footnotesize\setlength{\tabcolsep}{3pt}\renewcommand{\arraystretch}{1.25}
\begin{longtable}{p{2.13cm}p{2.13cm}p{2.13cm}p{2.13cm}p{2.13cm}p{2.13cm}p{2.13cm}}
\toprule
Оценка & Порог & Честных окон & Окон основной атаки & FPR & Pd & Ложных окон/ч \\
\midrule\endhead
nrf\_only & 9.00 & 304 & 272 & 0.214 & 0.871 & 76.97 \\
scp\_only & 6.91 & 304 & 272 & 0.257 & 0.496 & 92.37 \\
residual & 0.00 & 304 & 272 & 0.026 & 1.000 & 9.47 \\
residual\_quantized & 0.00 & 304 & 272 & 0.000 & 1.000 & 0.00 \\
nrf\_scp & 0.00 & 304 & 272 & 0.026 & 1.000 & 9.47 \\
full\_peer & 1.00 & 304 & 272 & 0.026 & 1.000 & 9.47 \\
\bottomrule\end{longtable}}

Окна в одной трассе зависимы. Значение FPR=0 означает отсутствие наблюдавшихся ложных окон в данном объёме, а не доказанный нулевой риск. Результаты при разных достигнутых FPR нельзя объявлять сравнением при одинаковом FPR=1\%.

Разделение по политике при общей SWRR-калибровке:

{\footnotesize\setlength{\tabcolsep}{3pt}\renewcommand{\arraystretch}{1.25}
\begin{longtable}{p{2.53cm}p{2.53cm}p{2.53cm}p{2.53cm}p{2.53cm}p{2.53cm}}
\toprule
Политика & Оценка & Честных окон & Окон атаки & FPR & Pd \\
\midrule\endhead
swrr & nrf\_only & 160 & 144 & 0.106 & 0.896 \\
swrr & scp\_only & 160 & 144 & 0.019 & 0.424 \\
swrr & residual & 160 & 144 & 0.050 & 1.000 \\
swrr & residual\_quantized & 160 & 144 & 0.000 & 1.000 \\
swrr & nrf\_scp & 160 & 144 & 0.050 & 1.000 \\
swrr & full\_peer & 160 & 144 & 0.050 & 1.000 \\
random & nrf\_only & 144 & 128 & 0.333 & 0.844 \\
random & scp\_only & 144 & 128 & 0.521 & 0.578 \\
random & residual & 144 & 128 & 0.000 & 1.000 \\
random & residual\_quantized & 144 & 128 & 0.000 & 1.000 \\
random & nrf\_scp & 144 & 128 & 0.000 & 1.000 \\
random & full\_peer & 144 & 128 & 0.000 & 1.000 \\
\bottomrule\end{longtable}}

\subsection*{Дополнительные кампании}

{\footnotesize\setlength{\tabcolsep}{3pt}\renewcommand{\arraystretch}{1.25}
\begin{longtable}{p{3.1cm}p{3.1cm}p{3.1cm}p{3.1cm}p{3.1cm}}
\toprule
Кампания & Доля SMF3, \% & \ensuremath{\Delta} к контролю, п.п. & Тревога residual & Тревога quantized \\
\midrule\endhead
adaptive-b0-parallel1 & 34.50 & -1.00 & 0.000 & 0.000 \\
adaptive-b1-parallel1 & 33.50 & -2.00 & 1.000 & 1.000 \\
adaptive-b2-parallel1 & 33.00 & -2.50 & 1.000 & 1.000 \\
adaptive-b4-parallel1 & 36.50 & 1.00 & 1.000 & 1.000 \\
adaptive-b8-parallel1 & 36.00 & 0.50 & 1.000 & 1.000 \\
adaptive-capacity150-parallel1 & 39.00 & 3.50 & 1.000 & 1.000 \\
adaptive-counterfactual-parallel1 & 35.50 & 0.00 & 0.000 & 0.000 \\
capacity150-parallel1 & 37.50 & 2.00 & 0.000 & 0.000 \\
honest-capacity70-parallel1 & 27.50 & — & 1.000 & 0.000 \\
onoff-parallel1 & 37.00 & 1.50 & 0.750 & 0.750 \\
ramp-parallel1 & 40.50 & 5.00 & 1.000 & 1.000 \\
\bottomrule\end{longtable}}

Доля тревог в кампании с включением/выключением не равна Pd по истинным интервалам активного искажения. Бюджет b=0 задаёт нулевое искажение. Пропущенное \ensuremath{\Delta} означает отсутствие подходящего парного контроля.

\subsection*{Распространение и свежесть отчётов}

{\footnotesize\setlength{\tabcolsep}{3pt}\renewcommand{\arraystretch}{1.25}
\begin{longtable}{p{3.1cm}p{3.1cm}p{3.1cm}p{3.1cm}p{3.1cm}}
\toprule
Величина & Событий & Медиана медиан, мс & Диапазон p95 прогонов, мс & Диапазон p99, мс \\
\midrule\endhead
Вычисление \ensuremath{\to} NRF & 38714 & 0.468 & 0.433 \ldots{} 0.965 & 0.517 \ldots{} 8.614 \\
Возраст при выборе AMF & 32340 & 518.434 & 924.353 \ldots{} 967.702 & 972.500 \ldots{} 997.531 \\
\bottomrule\end{longtable}}

Квантили отдельных прогонов не усредняются с выдачей за общий квантиль. Полные значения по каждому запуску сохранены в audit.json.

\subsection*{Проверяемый вывод}

Реализована и проверена цепочка динамического отчёта SMF \ensuremath{\to} NRF \ensuremath{\to} AMF и независимый учёт SCP. Величина и устойчивость смещения определяются только завершёнными парными опытами в таблице. Нельзя переносить вывод на штатную политику Open5GS, неизвестную ёмкость или неполное наблюдение SCP без новых опытов.

На завершённых дополнительных опытах простой residual выявляет искажение нагрузки, но не проверяет относительную capacity: при её завышении до 150 нагрузка может сообщаться честно, а распределение назначений всё равно меняется. Проверка capacity относительно доверенной конфигурации закрывает именно этот вариант. Честная ёмкость 70 выявляет другую проблему: сравнение целочисленного отчёта с дробным наблюдением даёт ложные тревоги. Вариант quantized устраняет эту причину в наблюдавшемся контроле; это не доказательство отсутствия иных причин ложных тревог.

Для малых бюджетов искажения эффект в отдельных коротких опытах может иметь отрицательный знак. Даже нулевое воздействие b=0 отличается от отдельного честного контрфактического прогона из-за конечного объёма и динамики выбора. Поэтому имеющиеся точки не устанавливают надёжную границу скрытности и полезности адаптивной атаки. Нужны более длинные серии с большим числом независимых повторов; ML на этих данных пока не обоснован.

\section*{Доступность материалов}

Полные патчи находятся в \texttt{patches/}, закреплённые версии — в \texttt{patches/manifest.json}, рецепт сборки — в \texttt{Dockerfile} и \texttt{compose.yaml}. Сырые журналы, конфигурации, PCAP, контрольные метрики, история операций UE и параметры опытов сохраняются в \texttt{runs/}. Таблицы анализа — в \texttt{results/article-evaluation/}; аналитический расчёт — отдельно в \texttt{results/theory/}. Проверка: \texttt{python scripts/test\_analysis.py}; обновление таблиц: \texttt{python scripts/evaluate\_article.py}; сборка этого текста: \texttt{python scripts/report\_article.py}.
\end{document}
